\documentclass[12pt,a4paper]{article}
\usepackage[utf8]{inputenc}
\usepackage[T1]{fontenc}
\usepackage{lmodern}
\usepackage[margin=1in]{geometry}
\usepackage{enumitem}
\usepackage{fancyhdr}
\usepackage{xcolor}
\usepackage{tcolorbox}
\usepackage{tabularx}

% Header/Footer
\pagestyle{fancy}
\fancyhf{}
\rhead{React + TypeScript Assessment}
\lhead{100 Questions}
\cfoot{\thepage}

\definecolor{easy}{rgb}{0.2,0.6,0.2}
\definecolor{medium}{rgb}{0.8,0.6,0.0}
\definecolor{hard}{rgb}{0.8,0.2,0.2}

\title{\Huge\textbf{React \& TypeScript Mastery Quiz}\\[0.5cm]
\Large 100 Multiple Choice Questions\\[0.3cm]
\normalsize Progressive Difficulty: Easy → Expert}
\author{Self-Assessment for First Principles Learning}
\date{January 2026}

\begin{document}

\maketitle
\thispagestyle{empty}

\vspace{1cm}

\begin{center}
\begin{tcolorbox}[width=0.8\textwidth,colback=gray!5,colframe=gray!50]
\textbf{Instructions:}
\begin{itemize}
    \item Questions 1-30: \textcolor{easy}{\textbf{Easy}} - Fundamental concepts
    \item Questions 31-60: \textcolor{medium}{\textbf{Medium}} - Applied knowledge
    \item Questions 61-85: \textcolor{hard}{\textbf{Hard}} - Advanced concepts
    \item Questions 86-100: \textcolor{hard}{\textbf{Expert}} - Deep understanding required
\end{itemize}
Each question has exactly ONE correct answer. Time yourself: 90 minutes recommended.
\end{tcolorbox}
\end{center}

\newpage

%==============================================================================
% EASY QUESTIONS (1-30)
%==============================================================================
\section*{\textcolor{easy}{Section 1: Easy (Questions 1-30)}}

\begin{enumerate}

% Q1 - Answer: C
\item What does HTML stand for?
\begin{enumerate}[label=\Alph*)]
    \item Hyper Text Making Language
    \item High Technical Modern Language
    \item Hyper Text Markup Language
    \item Home Tool Markup Language
\end{enumerate}

% Q2 - Answer: B
\item Which technology is used to style web pages?
\begin{enumerate}[label=\Alph*)]
    \item HTML
    \item CSS
    \item JavaScript
    \item TypeScript
\end{enumerate}

% Q3 - Answer: A
\item What is Node.js?
\begin{enumerate}[label=\Alph*)]
    \item A runtime environment that allows JavaScript to run outside browsers
    \item A web browser for developers
    \item A CSS framework
    \item A database management system
\end{enumerate}

% Q4 - Answer: D
\item What does npm stand for?
\begin{enumerate}[label=\Alph*)]
    \item New Programming Method
    \item Node Project Manager
    \item Network Package Module
    \item Node Package Manager
\end{enumerate}

% Q5 - Answer: B
\item Which command creates a new Vite project?
\begin{enumerate}[label=\Alph*)]
    \item \texttt{vite new project}
    \item \texttt{npm create vite@latest}
    \item \texttt{create-vite-app}
    \item \texttt{npm init vite}
\end{enumerate}

% Q6 - Answer: C
\item What is the purpose of the \texttt{node\_modules} folder?
\begin{enumerate}[label=\Alph*)]
    \item Store your source code
    \item Store images and assets
    \item Store downloaded packages/dependencies
    \item Store compiled output
\end{enumerate}

% Q7 - Answer: A
\item Which file extension indicates a TypeScript React component?
\begin{enumerate}[label=\Alph*)]
    \item \texttt{.tsx}
    \item \texttt{.jsx}
    \item \texttt{.ts}
    \item \texttt{.react}
\end{enumerate}

% Q8 - Answer: D
\item What is JSX?
\begin{enumerate}[label=\Alph*)]
    \item A JavaScript testing framework
    \item A CSS preprocessor
    \item A database query language
    \item HTML-like syntax that compiles to JavaScript
\end{enumerate}

% Q9 - Answer: B
\item In React, what is a component?
\begin{enumerate}[label=\Alph*)]
    \item A CSS class
    \item A reusable piece of UI
    \item A JavaScript variable
    \item A HTML element
\end{enumerate}

% Q10 - Answer: C
\item Which command starts the Vite development server?
\begin{enumerate}[label=\Alph*)]
    \item \texttt{npm start}
    \item \texttt{npm server}
    \item \texttt{npm run dev}
    \item \texttt{npm run start}
\end{enumerate}

% Q11 - Answer: A
\item What does ``localhost'' refer to?
\begin{enumerate}[label=\Alph*)]
    \item Your own computer
    \item A remote server
    \item The internet
    \item A cloud service
\end{enumerate}

% Q12 - Answer: D
\item What is Git?
\begin{enumerate}[label=\Alph*)]
    \item A code editor
    \item A programming language
    \item A web browser
    \item A version control system
\end{enumerate}

% Q13 - Answer: B
\item Which command initializes a new Git repository?
\begin{enumerate}[label=\Alph*)]
    \item \texttt{git start}
    \item \texttt{git init}
    \item \texttt{git new}
    \item \texttt{git create}
\end{enumerate}

% Q14 - Answer: C
\item What does \texttt{git add .} do?
\begin{enumerate}[label=\Alph*)]
    \item Creates a new file
    \item Removes all files
    \item Stages all changes for commit
    \item Pushes to remote repository
\end{enumerate}

% Q15 - Answer: A
\item In TypeScript, how do you declare a variable that holds a string?
\begin{enumerate}[label=\Alph*)]
    \item \texttt{let name: string = "Sofia";}
    \item \texttt{let name = string("Sofia");}
    \item \texttt{string name = "Sofia";}
    \item \texttt{var name as string = "Sofia";}
\end{enumerate}

% Q16 - Answer: D
\item What is the purpose of \texttt{export default}?
\begin{enumerate}[label=\Alph*)]
    \item Delete a component
    \item Import a component
    \item Style a component
    \item Make a component available to other files
\end{enumerate}

% Q17 - Answer: B
\item Which HTML tag represents the main content of a document?
\begin{enumerate}[label=\Alph*)]
    \item \texttt{<content>}
    \item \texttt{<main>}
    \item \texttt{<body>}
    \item \texttt{<div>}
\end{enumerate}

% Q18 - Answer: C
\item In CSS, what does \texttt{margin: 0 auto;} do?
\begin{enumerate}[label=\Alph*)]
    \item Removes all margins
    \item Adds automatic margins everywhere
    \item Centers the element horizontally
    \item Centers the element vertically
\end{enumerate}

% Q19 - Answer: A
\item What is the correct syntax for a React function component?
\begin{enumerate}[label=\Alph*)]
    \item \texttt{function MyComponent() \{ return <div>Hello</div>; \}}
    \item \texttt{class MyComponent \{ render <div>Hello</div> \}}
    \item \texttt{component MyComponent() \{ <div>Hello</div> \}}
    \item \texttt{def MyComponent(): return <div>Hello</div>}
\end{enumerate}

% Q20 - Answer: D
\item Which symbol is used for single-line comments in JavaScript?
\begin{enumerate}[label=\Alph*)]
    \item \texttt{\#}
    \item \texttt{<!-- -->}
    \item \texttt{/* */}
    \item \texttt{//}
\end{enumerate}

% Q21 - Answer: B
\item What does CSS stand for?
\begin{enumerate}[label=\Alph*)]
    \item Computer Style System
    \item Cascading Style Sheets
    \item Creative Styling Standard
    \item Coded Style Syntax
\end{enumerate}

% Q22 - Answer: C
\item In React, how do you render a variable inside JSX?
\begin{enumerate}[label=\Alph*)]
    \item \texttt{(variable)}
    \item \texttt{[variable]}
    \item \texttt{\{variable\}}
    \item \texttt{<variable>}
\end{enumerate}

% Q23 - Answer: A
\item What is the difference between \texttt{==} and \texttt{===} in JavaScript?
\begin{enumerate}[label=\Alph*)]
    \item \texttt{===} checks type and value, \texttt{==} only checks value
    \item They are identical
    \item \texttt{==} is for numbers, \texttt{===} is for strings
    \item \texttt{===} is deprecated
\end{enumerate}

% Q24 - Answer: D
\item Which command pushes commits to a remote repository?
\begin{enumerate}[label=\Alph*)]
    \item \texttt{git upload}
    \item \texttt{git send}
    \item \texttt{git commit}
    \item \texttt{git push}
\end{enumerate}

% Q25 - Answer: B
\item What is the purpose of the \texttt{.gitignore} file?
\begin{enumerate}[label=\Alph*)]
    \item List files to delete
    \item List files Git should not track
    \item List important files
    \item Configure Git settings
\end{enumerate}

% Q26 - Answer: C
\item In TypeScript, what does \texttt{?} after a property name indicate?
\begin{enumerate}[label=\Alph*)]
    \item The property is required
    \item The property is deprecated
    \item The property is optional
    \item The property is read-only
\end{enumerate}

% Q27 - Answer: A
\item What is the purpose of \texttt{index.html} in a React project?
\begin{enumerate}[label=\Alph*)]
    \item It's the HTML shell where React mounts
    \item It contains all CSS styles
    \item It stores component data
    \item It configures the build process
\end{enumerate}

% Q28 - Answer: D
\item Which CSS property changes text color?
\begin{enumerate}[label=\Alph*)]
    \item \texttt{text-color}
    \item \texttt{font-color}
    \item \texttt{background-color}
    \item \texttt{color}
\end{enumerate}

% Q29 - Answer: B
\item What does the \texttt{cd} command do in terminal?
\begin{enumerate}[label=\Alph*)]
    \item Copy directory
    \item Change directory
    \item Create directory
    \item Clear directory
\end{enumerate}

% Q30 - Answer: C
\item In a React component, what does the \texttt{return} statement contain?
\begin{enumerate}[label=\Alph*)]
    \item CSS styles
    \item JavaScript logic only
    \item JSX/HTML to be rendered
    \item Data to be saved
\end{enumerate}

%==============================================================================
% MEDIUM QUESTIONS (31-60)
%==============================================================================
\newpage
\section*{\textcolor{medium}{Section 2: Medium (Questions 31-60)}}

% Q31 - Answer: B
\item What is destructuring in JavaScript?
\begin{enumerate}[label=\Alph*)]
    \item Deleting objects from memory
    \item Extracting values from objects/arrays into variables
    \item Breaking code into smaller files
    \item Converting objects to strings
\end{enumerate}

% Q32 - Answer: D
\item Given \texttt{const [count, setCount] = useState(0)}, what does \texttt{setCount} do?
\begin{enumerate}[label=\Alph*)]
    \item Gets the current count
    \item Resets count to zero
    \item Deletes the count variable
    \item Updates the count value and triggers re-render
\end{enumerate}

% Q33 - Answer: A
\item What is the purpose of \texttt{key} prop when rendering lists in React?
\begin{enumerate}[label=\Alph*)]
    \item Helps React identify which items changed
    \item Encrypts the list data
    \item Styles the list items
    \item Sorts the list automatically
\end{enumerate}

% Q34 - Answer: C
\item What does \texttt{Array.map()} return?
\begin{enumerate}[label=\Alph*)]
    \item A boolean
    \item The original array modified
    \item A new array with transformed elements
    \item Nothing (undefined)
\end{enumerate}

% Q35 - Answer: B
\item In \texttt{const \{ name, age \} = person}, what is this syntax called?
\begin{enumerate}[label=\Alph*)]
    \item Spread operator
    \item Object destructuring
    \item Rest parameters
    \item Template literals
\end{enumerate}

% Q36 - Answer: D
\item What is the difference between \texttt{props} and \texttt{state} in React?
\begin{enumerate}[label=\Alph*)]
    \item Props are mutable, state is immutable
    \item They are the same thing
    \item Props are for styling, state is for logic
    \item Props come from parent, state is managed internally
\end{enumerate}

% Q37 - Answer: A
\item What does \texttt{npm install} do?
\begin{enumerate}[label=\Alph*)]
    \item Downloads dependencies listed in package.json
    \item Creates a new npm project
    \item Uploads your project to npm
    \item Removes all packages
\end{enumerate}

% Q38 - Answer: C
\item In TypeScript, what is an interface used for?
\begin{enumerate}[label=\Alph*)]
    \item Creating visual components
    \item Defining CSS styles
    \item Defining the shape/structure of objects
    \item Connecting to databases
\end{enumerate}

% Q39 - Answer: B
\item What does the \texttt{children} prop represent in React?
\begin{enumerate}[label=\Alph*)]
    \item Child components in a separate file
    \item Content placed between component's opening and closing tags
    \item Nested CSS rules
    \item Sub-routes in routing
\end{enumerate}

% Q40 - Answer: D
\item What is the purpose of \texttt{useEffect} hook?
\begin{enumerate}[label=\Alph*)]
    \item Create visual effects
    \item Define component styles
    \item Handle form inputs
    \item Perform side effects (API calls, DOM updates, etc.)
\end{enumerate}

% Q41 - Answer: A
\item What does CSS \texttt{display: flex} enable?
\begin{enumerate}[label=\Alph*)]
    \item Flexbox layout
    \item Grid layout
    \item Float layout
    \item Absolute positioning
\end{enumerate}

% Q42 - Answer: C
\item In \texttt{useEffect(() => \{\}, [dependency])}, what does the dependency array do?
\begin{enumerate}[label=\Alph*)]
    \item Lists functions to call
    \item Lists components to render
    \item Controls when the effect re-runs
    \item Lists variables to export
\end{enumerate}

% Q43 - Answer: B
\item What does \texttt{event.preventDefault()} do?
\begin{enumerate}[label=\Alph*)]
    \item Prevents rendering
    \item Stops the default browser behavior for that event
    \item Prevents state updates
    \item Stops event bubbling
\end{enumerate}

% Q44 - Answer: D
\item What is the correct way to conditionally render in JSX?
\begin{enumerate}[label=\Alph*)]
    \item \texttt{if (condition) <Component />}
    \item \texttt{<if condition><Component /></if>}
    \item \texttt{<Component if=\{condition\} />}
    \item \texttt{\{condition \&\& <Component />\}}
\end{enumerate}

% Q45 - Answer: A
\item In CSS, what is the BEM naming convention?
\begin{enumerate}[label=\Alph*)]
    \item Block\_\_Element--Modifier
    \item Basic Element Method
    \item Browser Extension Module
    \item Build Element Modify
\end{enumerate}

% Q46 - Answer: C
\item What does \texttt{:root} select in CSS?
\begin{enumerate}[label=\Alph*)]
    \item The body element
    \item The first div
    \item The html element (document root)
    \item All elements
\end{enumerate}

% Q47 - Answer: B
\item What is a TypeScript union type?
\begin{enumerate}[label=\Alph*)]
    \item Combining two objects
    \item A type that can be one of several types (e.g., \texttt{string | number})
    \item Merging two arrays
    \item Joining two functions
\end{enumerate}

% Q48 - Answer: D
\item What does \texttt{localStorage.setItem('key', 'value')} do?
\begin{enumerate}[label=\Alph*)]
    \item Deletes data from browser
    \item Reads data from a file
    \item Sends data to a server
    \item Stores data in the browser that persists after closing
\end{enumerate}

% Q49 - Answer: A
\item What is the purpose of React Context?
\begin{enumerate}[label=\Alph*)]
    \item Share data across components without prop drilling
    \item Define component context/description
    \item Handle routing
    \item Manage CSS styles
\end{enumerate}

% Q50 - Answer: C
\item What does \texttt{flex-direction: column} do?
\begin{enumerate}[label=\Alph*)]
    \item Arranges items horizontally
    \item Hides items
    \item Arranges items vertically (stacked)
    \item Rotates items 90 degrees
\end{enumerate}

% Q51 - Answer: B
\item In TypeScript, what does \texttt{: void} indicate for a function?
\begin{enumerate}[label=\Alph*)]
    \item The function takes no parameters
    \item The function returns nothing
    \item The function is empty
    \item The function is deprecated
\end{enumerate}

% Q52 - Answer: D
\item What is the purpose of \texttt{useRef} hook?
\begin{enumerate}[label=\Alph*)]
    \item Create references to other files
    \item Manage state that triggers re-renders
    \item Handle API requests
    \item Reference DOM elements or store mutable values without re-render
\end{enumerate}

% Q53 - Answer: A
\item What does \texttt{git commit -m "message"} do?
\begin{enumerate}[label=\Alph*)]
    \item Creates a snapshot of staged changes with a description
    \item Sends changes to remote server
    \item Deletes previous commits
    \item Merges branches
\end{enumerate}

% Q54 - Answer: C
\item In CSS, what does \texttt{transition: color 0.3s ease} do?
\begin{enumerate}[label=\Alph*)]
    \item Changes color instantly
    \item Removes color
    \item Animates color change over 0.3 seconds
    \item Delays color change by 0.3 seconds
\end{enumerate}

% Q55 - Answer: B
\item What does the spread operator (\texttt{...}) do?
\begin{enumerate}[label=\Alph*)]
    \item Multiplies values
    \item Expands an array/object into individual elements
    \item Creates a range of numbers
    \item Concatenates strings
\end{enumerate}

% Q56 - Answer: D
\item What is the purpose of \texttt{package.json}?
\begin{enumerate}[label=\Alph*)]
    \item Store application data
    \item Configure Git settings
    \item Define HTML structure
    \item Define project metadata and dependencies
\end{enumerate}

% Q57 - Answer: A
\item What does \texttt{target="\_blank"} do on an anchor tag?
\begin{enumerate}[label=\Alph*)]
    \item Opens link in a new tab
    \item Opens link in same tab
    \item Downloads the link
    \item Hides the link
\end{enumerate}

% Q58 - Answer: C
\item In TypeScript, what does \texttt{Experience[]} mean?
\begin{enumerate}[label=\Alph*)]
    \item An Experience object with an array inside
    \item The length of Experience
    \item An array of Experience objects
    \item An optional Experience
\end{enumerate}

% Q59 - Answer: B
\item What does CSS \texttt{gap: 1rem} do in a flex container?
\begin{enumerate}[label=\Alph*)]
    \item Creates space outside the container
    \item Creates space between flex items
    \item Creates space inside each item
    \item Removes all spacing
\end{enumerate}

% Q60 - Answer: D
\item What is the purpose of \texttt{aria-label} attribute?
\begin{enumerate}[label=\Alph*)]
    \item Style the element
    \item Add a visible label
    \item Create an animation
    \item Provide text for screen readers (accessibility)
\end{enumerate}

%==============================================================================
% HARD QUESTIONS (61-85)
%==============================================================================
\newpage
\section*{\textcolor{hard}{Section 3: Hard (Questions 61-85)}}

% Q61 - Answer: C
\item In the code \texttt{useState(() => computeInitialValue())}, why use a function?
\begin{enumerate}[label=\Alph*)]
    \item It's required syntax
    \item To make the code shorter
    \item To defer expensive computation until actually needed (lazy initialization)
    \item To make the value immutable
\end{enumerate}

% Q62 - Answer: A
\item What happens when you call \texttt{setState} with the same value as current state?
\begin{enumerate}[label=\Alph*)]
    \item React bails out and doesn't re-render
    \item React always re-renders
    \item An error is thrown
    \item The component unmounts
\end{enumerate}

% Q63 - Answer: D
\item In \texttt{useEffect}, what does returning a function do?
\begin{enumerate}[label=\Alph*)]
    \item Prevents the effect from running
    \item Caches the effect result
    \item Delays the effect
    \item Provides cleanup when component unmounts or before re-run
\end{enumerate}

% Q64 - Answer: B
\item What is the issue with \texttt{onClick=\{handleClick()\}} vs \texttt{onClick=\{handleClick\}}?
\begin{enumerate}[label=\Alph*)]
    \item No difference
    \item With (), the function executes immediately on render instead of on click
    \item Without (), the function never runs
    \item With () is the correct syntax
\end{enumerate}

% Q65 - Answer: C
\item What is ``prop drilling'' in React?
\begin{enumerate}[label=\Alph*)]
    \item Testing props in components
    \item Removing unused props
    \item Passing props through many levels of components
    \item Creating props with default values
\end{enumerate}

% Q66 - Answer: A
\item Why might you use \texttt{useCallback} hook?
\begin{enumerate}[label=\Alph*)]
    \item To memoize a function and prevent unnecessary re-creation
    \item To call a function after render
    \item To handle async callbacks
    \item To create global callbacks
\end{enumerate}

% Q67 - Answer: D
\item In TypeScript, what does \texttt{as const} do?
\begin{enumerate}[label=\Alph*)]
    \item Converts to a constant variable
    \item Creates a copy
    \item Enables strict mode
    \item Makes the type literal and readonly (const assertion)
\end{enumerate}

% Q68 - Answer: B
\item What is the purpose of \texttt{React.StrictMode}?
\begin{enumerate}[label=\Alph*)]
    \item Enforces TypeScript strict mode
    \item Highlights potential problems by running effects twice in dev
    \item Prevents any errors from occurring
    \item Enables production optimizations
\end{enumerate}

% Q69 - Answer: C
\item What does \texttt{document.getElementById('root')!} mean in TypeScript?
\begin{enumerate}[label=\Alph*)]
    \item Negates the result
    \item Checks if element exists
    \item Non-null assertion (tells TypeScript it won't be null)
    \item Makes the call optional
\end{enumerate}

% Q70 - Answer: A
\item What is the difference between \texttt{createContext(undefined)} and \texttt{createContext(null)}?
\begin{enumerate}[label=\Alph*)]
    \item \texttt{undefined} allows detecting missing provider, \texttt{null} is a valid value
    \item They are identical
    \item \texttt{null} is preferred for performance
    \item \texttt{undefined} causes errors
\end{enumerate}

% Q71 - Answer: D
\item In CSS, what does \texttt{preserve\-Aspect\-Ratio="xMidYMid meet"} do on SVG?
\begin{enumerate}[label=\Alph*)]
    \item Stretches SVG to fill container
    \item Crops the SVG
    \item Rotates the SVG
    \item Scales uniformly while centering and fitting within container
\end{enumerate}

% Q72 - Answer: B
\item What is the CSS specificity order (lowest to highest)?
\begin{enumerate}[label=\Alph*)]
    \item ID, Class, Element
    \item Element, Class, ID
    \item Class, ID, Element
    \item Element, ID, Class
\end{enumerate}

% Q73 - Answer: C
\item Why use \texttt{rel="noopener noreferrer"} with \texttt{target="\_blank"}?
\begin{enumerate}[label=\Alph*)]
    \item For SEO purposes
    \item To improve loading speed
    \item To prevent the new page from accessing the original page (security)
    \item For accessibility
\end{enumerate}

% Q74 - Answer: A
\item What is a closure in JavaScript?
\begin{enumerate}[label=\Alph*)]
    \item A function that has access to variables from its outer scope
    \item A way to close browser windows
    \item A method to end loops
    \item A type of object
\end{enumerate}

% Q75 - Answer: D
\item What does \texttt{npm run build} typically produce?
\begin{enumerate}[label=\Alph*)]
    \item A new npm package
    \item Source maps only
    \item Node modules
    \item Optimized, minified production files in dist folder
\end{enumerate}

% Q76 - Answer: B
\item In \texttt{setCount(prev => prev + 1)}, why use a function instead of \texttt{setCount(count + 1)}?
\begin{enumerate}[label=\Alph*)]
    \item It's faster
    \item Ensures correct value when multiple updates are batched
    \item It's required syntax
    \item For TypeScript compatibility
\end{enumerate}

% Q77 - Answer: C
\item What is the virtual DOM?
\begin{enumerate}[label=\Alph*)]
    \item A backup of the real DOM
    \item The DOM in virtual reality
    \item A lightweight JavaScript representation of the DOM for efficient updates
    \item The DOM stored in memory cache
\end{enumerate}

% Q78 - Answer: A
\item What does \texttt{Object.freeze(obj)} do?
\begin{enumerate}[label=\Alph*)]
    \item Makes the object immutable (can't add/modify/delete properties)
    \item Stores object in cache
    \item Creates a copy of the object
    \item Converts object to string
\end{enumerate}

% Q79 - Answer: D
\item Why might CSS custom properties (variables) be preferred over preprocessor variables?
\begin{enumerate}[label=\Alph*)]
    \item They compile faster
    \item They use less memory
    \item They're easier to write
    \item They can be changed at runtime and scoped to elements
\end{enumerate}

% Q80 - Answer: B
\item In TypeScript, what is a generic type?
\begin{enumerate}[label=\Alph*)]
    \item A type that works everywhere
    \item A type parameter that makes functions/classes work with multiple types
    \item A default type
    \item A type with no restrictions
\end{enumerate}

% Q81 - Answer: C
\item What is reconciliation in React?
\begin{enumerate}[label=\Alph*)]
    \item Combining multiple components
    \item Resolving merge conflicts
    \item The algorithm React uses to diff the virtual DOM and update real DOM
    \item Connecting to a database
\end{enumerate}

% Q82 - Answer: A
\item What does \texttt{window.matchMedia('(prefers-color-scheme: dark)')} return?
\begin{enumerate}[label=\Alph*)]
    \item A MediaQueryList indicating if user prefers dark mode
    \item A boolean true/false
    \item The current color scheme
    \item A CSS style object
\end{enumerate}

% Q83 - Answer: D
\item Why is it important to keep the key prop stable across re-renders?
\begin{enumerate}[label=\Alph*)]
    \item For security
    \item For SEO
    \item For styling
    \item To preserve component state and avoid unnecessary DOM operations
\end{enumerate}

% Q84 - Answer: B
\item What is the difference between \texttt{npm} and \texttt{npx}?
\begin{enumerate}[label=\Alph*)]
    \item npx is faster
    \item npx executes packages without installing, npm manages packages
    \item npm is deprecated
    \item They're the same
\end{enumerate}

% Q85 - Answer: C
\item In CSS, what does \texttt{box-sizing: border-box} change?
\begin{enumerate}[label=\Alph*)]
    \item Adds a border to all boxes
    \item Removes padding from elements
    \item Makes width/height include padding and border, not just content
    \item Centers the element
\end{enumerate}

%==============================================================================
% EXPERT QUESTIONS (86-100)
%==============================================================================
\newpage
\section*{\textcolor{hard}{Section 4: Expert (Questions 86-100)}}

% Q86 - Answer: D
\item Consider this code:
\begin{verbatim}
const [items, setItems] = useState([1, 2, 3]);
const addItem = () => {
  items.push(4);
  setItems(items);
};
\end{verbatim}
Why won't clicking the button trigger a re-render?
\begin{enumerate}[label=\Alph*)]
    \item \texttt{useState} doesn't work with arrays
    \item \texttt{push} is not a valid array method
    \item The syntax is incorrect
    \item Same array reference is passed; React uses referential equality
\end{enumerate}

% Q87 - Answer: A
\item What is the time complexity of React's reconciliation algorithm?
\begin{enumerate}[label=\Alph*)]
    \item O(n) where n is number of elements
    \item O(n\textsuperscript{2})
    \item O(n\textsuperscript{3})
    \item O(log n)
\end{enumerate}

% Q88 - Answer: C
\item In this TypeScript code, what is the type of \texttt{result}?
\begin{verbatim}
const obj = { a: 1, b: "hello" } as const;
type Result = typeof obj;
\end{verbatim}
\begin{enumerate}[label=\Alph*)]
    \item \texttt{\{ a: number, b: string \}}
    \item \texttt{\{ a: any, b: any \}}
    \item \texttt{\{ readonly a: 1, readonly b: "hello" \}}
    \item \texttt{Object}
\end{enumerate}

% Q89 - Answer: B
\item What is the stale closure problem in React hooks?
\begin{enumerate}[label=\Alph*)]
    \item Using old hook syntax
    \item Callbacks capturing outdated values from previous renders
    \item Memory leaks from closures
    \item Infinite loops in useEffect
\end{enumerate}

% Q90 - Answer: D
\item Why might you wrap a context value in \texttt{useMemo}?
\begin{enumerate}[label=\Alph*)]
    \item To make the value immutable
    \item For TypeScript compatibility
    \item To enable caching
    \item To prevent unnecessary re-renders of all consumers on provider re-render
\end{enumerate}

% Q91 - Answer: A
\item What is tree shaking in JavaScript bundlers?
\begin{enumerate}[label=\Alph*)]
    \item Removing unused code from the final bundle
    \item Reorganizing the component tree
    \item Optimizing recursive functions
    \item Creating backup copies
\end{enumerate}

% Q92 - Answer: C
\item In TypeScript, what does this type represent?
\begin{verbatim}
type NonNullable<T> = T extends null | undefined ? never : T;
\end{verbatim}
\begin{enumerate}[label=\Alph*)]
    \item A type that only allows null
    \item A generic parameter
    \item A conditional type that excludes null and undefined
    \item A type guard
\end{enumerate}

% Q93 - Answer: B
\item What is the purpose of React Fiber?
\begin{enumerate}[label=\Alph*)]
    \item Network optimization
    \item Incremental rendering with the ability to pause, abort, or reuse work
    \item CSS-in-JS implementation
    \item Server-side rendering
\end{enumerate}

% Q94 - Answer: D
\item Given this CSS:
\begin{verbatim}
:root { --color: blue; }
.dark { --color: white; }
.component { color: var(--color); }
\end{verbatim}
If \texttt{<div class="dark"><span class="component">Text</span></div>}, what color is Text?
\begin{enumerate}[label=\Alph*)]
    \item Blue (root always wins)
    \item Undefined (error)
    \item Black (default)
    \item White (inherits from parent's custom property)
\end{enumerate}

% Q95 - Answer: A
\item What is the difference between controlled and uncontrolled components?
\begin{enumerate}[label=\Alph*)]
    \item Controlled: React state drives the value; Uncontrolled: DOM holds the value
    \item Controlled: read-only; Uncontrolled: editable
    \item They're the same
    \item Controlled is for forms, uncontrolled is for display
\end{enumerate}

% Q96 - Answer: C
\item In this code, what's the problem?
\begin{verbatim}
useEffect(() => {
  const data = await fetchData();
  setData(data);
}, []);
\end{verbatim}
\begin{enumerate}[label=\Alph*)]
    \item fetchData might return undefined
    \item setData is called incorrectly
    \item useEffect callback cannot be async; await is invalid here
    \item The dependency array is wrong
\end{enumerate}

% Q97 - Answer: B
\item What is the render-commit phase distinction in React?
\begin{enumerate}[label=\Alph*)]
    \item Render is server-side, commit is client-side
    \item Render calculates changes (can be paused), commit applies them to DOM (cannot)
    \item Render is development, commit is production
    \item They're the same phase
\end{enumerate}

% Q98 - Answer: D
\item In TypeScript, what does this mapped type do?
\begin{verbatim}
type Readonly<T> = { readonly [P in keyof T]: T[P] };
\end{verbatim}
\begin{enumerate}[label=\Alph*)]
    \item Removes all properties
    \item Makes all properties optional
    \item Creates a copy of T
    \item Makes all properties readonly by iterating over keys of T
\end{enumerate}

% Q99 - Answer: A
\item What is the purpose of \texttt{React.memo()}?
\begin{enumerate}[label=\Alph*)]
    \item Memoizes a component to prevent re-render if props haven't changed
    \item Stores component in memory cache
    \item Creates a memo/note about the component
    \item Enables strict mode for a component
\end{enumerate}

% Q100 - Answer: C
\item Consider this scenario: A parent component re-renders. It passes an inline object \texttt{\{id: 1\}} as a prop to a child wrapped in \texttt{React.memo()}. Will the child re-render?
\begin{enumerate}[label=\Alph*)]
    \item No, the object values are the same
    \item No, React.memo prevents all re-renders
    \item Yes, because a new object reference is created each render
    \item Only if the child has state
\end{enumerate}

\end{enumerate}

%==============================================================================
% ANSWER KEY
%==============================================================================
\newpage
\section*{Answer Key}

\begin{center}
\textbf{Tear off or cover this page while taking the quiz.}
\end{center}

\vspace{1cm}

\begin{center}
\begin{tabular}{|c|c||c|c||c|c||c|c|}
\hline
\textbf{Q} & \textbf{Ans} & \textbf{Q} & \textbf{Ans} & \textbf{Q} & \textbf{Ans} & \textbf{Q} & \textbf{Ans} \\
\hline
1 & C & 26 & C & 51 & B & 76 & B \\
2 & B & 27 & A & 52 & D & 77 & C \\
3 & A & 28 & D & 53 & A & 78 & A \\
4 & D & 29 & B & 54 & C & 79 & D \\
5 & B & 30 & C & 55 & B & 80 & B \\
6 & C & 31 & B & 56 & D & 81 & C \\
7 & A & 32 & D & 57 & A & 82 & A \\
8 & D & 33 & A & 58 & C & 83 & D \\
9 & B & 34 & C & 59 & B & 84 & B \\
10 & C & 35 & B & 60 & D & 85 & C \\
11 & A & 36 & D & 61 & C & 86 & D \\
12 & D & 37 & A & 62 & A & 87 & A \\
13 & B & 38 & C & 63 & D & 88 & C \\
14 & C & 39 & B & 64 & B & 89 & B \\
15 & A & 40 & D & 65 & C & 90 & D \\
16 & D & 41 & A & 66 & A & 91 & A \\
17 & B & 42 & C & 67 & D & 92 & C \\
18 & C & 43 & B & 68 & B & 93 & B \\
19 & A & 44 & D & 69 & C & 94 & D \\
20 & D & 45 & A & 70 & A & 95 & A \\
21 & B & 46 & C & 71 & D & 96 & C \\
22 & C & 47 & B & 72 & B & 97 & B \\
23 & A & 48 & D & 73 & C & 98 & D \\
24 & D & 49 & A & 74 & A & 99 & A \\
25 & B & 50 & C & 75 & D & 100 & C \\
\hline
\end{tabular}
\end{center}

\vspace{1cm}

\subsection*{Scoring Guide}
\begin{itemize}
    \item \textbf{90-100}: Expert level - Ready for production React/TypeScript
    \item \textbf{75-89}: Advanced - Strong foundation, review missed concepts
    \item \textbf{60-74}: Intermediate - Good progress, more practice needed
    \item \textbf{45-59}: Beginner+ - Revisit tutorial chapters
    \item \textbf{Below 45}: Beginner - Start from Chapter 1 again
\end{itemize}

\subsection*{Answer Distribution}
\begin{tabular}{|c|c|}
\hline
\textbf{Answer} & \textbf{Count} \\
\hline
A & 25 \\
B & 25 \\
C & 26 \\
D & 24 \\
\hline
\end{tabular}

\vspace{0.5cm}
Answers are distributed roughly evenly and randomized---no patterns to exploit!

\end{document}
